\section{Benchmark}
\label{sec:benchmark}

\paragraph{Construction.}
We build the benchmark from six public IoT and smart-home corpora that appear in the
retrieval and long-context literature. For each source we assemble excerpts in which
records for several entities of the same type are interleaved, and we inject synthetic
protected values (account identifiers, contact addresses, phone numbers) into the records
that already exist, rather than into separate carriers. Reference answers are derived
deterministically from the injected structure, so no answer is authored by hand.

The injected values are canaries, not real personal data: public corpora cannot carry real
PII, and using a canary makes leakage measurable exactly. The consequence is that our
privacy claim concerns the recoverability of \emph{these} values, not of arbitrary
personal information that may co-occur in real logs.

\paragraph{Question form.}
Each case poses an entity-resolving question: the question names an account and asks for an
attribute recorded for it (which device is associated, which persona identifier is stored,
and so on). Two families of excerpts are included. In \emph{multi-entity} excerpts the same
question template applies to six accounts and the target account is distinguished only by
its identifier. In \emph{needle} excerpts the question names a specific record identifier
and asks for a field value inside it.

\paragraph{Sizes.}
We report three sets. The \emph{main} set contains $112$ cases, $20$ from each of five
sources and $12$ from a sixth, and is used for the main table, the budget curve, the
cross-model and the paraphrasing experiments. The \emph{extension} set contains $2{,}728$
cases from the same six sources and is used to check that the main-set conclusions survive
an order-of-magnitude increase in size and a different source mix. The six sources contribute
$426$, $600$, $312$, $490$, $300$ and $600$ cases respectively, so no single source exceeds
$22\%$ of the extension. Both sets share the same question templates.

A note on key hygiene: the two sets reuse the same case identifiers for different content,
and joining results across them by identifier silently mixes records. Every analysis
reported here (i) resolves results within a single set, and (ii) verifies the join by
comparing the uncompressed-context arm against a reference value computed from the same set.
The check is reported alongside each table.

\paragraph{Public-benchmark variants.}
To test the rule outside our own construction we derive three variants from LongBench
passages, $300$ cases each, with the same injected canaries. \emph{Topical} leaves the
original questions, which name no entity. \emph{Single-hop} replaces the question with one
that names an account, and places the answer in the same record as the account.
\emph{Two-hop} places the answer six sentences away from the account mention, so that
resolving the question requires following the record structure rather than the local
window. Section~\ref{sec:analysis} reports all three, including the one we do not claim.
